\cleardoublepage
\chapter{KODE PROGRAM}

\section{Repositori dan Strukturnya}
Kode program yang diimplementasikan untuk membangun sistem Equilibria disimpan dalam repositori daring yang dapat diakses pada pranala
\url{https://github.com/RunningPie/Equilibria}

Struktur repositori yang digunakan adalah seperti di bawah ini.

\dirtree{%
.1 equilibria/.
.2 LICENSE.
.2 README.md.
.2 client/.
.3 README.md.
.3 index.html.
.3 public/.
.3 src/.
.2 docs/.
.3 18222047-ProposalTA-signed.pdf.
.3 \parbox{11cm}{compressed Adaptive Assessment and Content Recommendation in Online Programming Courses.pdf}.
.2 server/.
.3 Dockerfile.
.3 alembic/.
.3 alembic.ini.
.3 app/.
.3 docker-compose.yml.
.3 logs/.
.3 requirements.txt.
}


\section{Implementasi Kode Teknis}
Beberapa algoritma yang dinilai penting untuk dilampirkan implementasi secara teknisnya ada dua, yaitu algoritma untuk menghitung Elo rating dan algoritma untuk memilih rekan sejawat (re-ranking).

\subsection{Algoritma Elo Rating}
Kode implementasi untuk Elo rating dapat ditinjau pada repositori dalam berkas bernama \texttt{elo\_engine.py} di direktori \texttt{app/core/}. Berikut adalah kode lengkapnya:

\begin{lstlisting}[language=Python, caption={Implementasi Algoritma Elo Rating}, label={lst:elo_rating}]
import math
import numpy
from typing import Tuple
from sqlalchemy.future import select
from sqlalchemy.ext.asyncio import AsyncSession
from app.db.models import AssessmentLog

# === Konstanta sesuai Vesin et al. (2022) ===
BASE_RATING = 1300.0  # Rating awal semua siswa
RATING_MIN = 1000.0
RATING_MAX = 1800.0  
K_FACTORS = {
    'novice': 30,    # 0-9 attempt
    'intermediate': 20,  # 10-24 attempt
    'advanced': 15,  # 25-49 attempt
    'expert': 10     # 50+ attempt
}
TIME_DISCRIMINATION = 1e-6 
DEFAULT_TIME_LIMIT = 300000  # 5 menit dalam milidetik

def calculate_initial_theta(correct_count: int, total_questions: int = 5) -> float:
    base_rating = 1300.0
    baseline_correct = 2.5
    multiplier = 80.0
    adjustment = (correct_count - baseline_correct) * multiplier
    theta = base_rating + adjustment
    return max(1100, min(1500, theta))

def calculate_expected_score(student_rating: float, question_difficulty: float)->float:
    rating_diff = question_difficulty - student_rating
    return 1.0 / (1.0 + math.pow(10, rating_diff / 400.0))

def get_k_factor(total_attempts: int)->int:
    if total_attempts < 10:
        return K_FACTORS['novice']
    elif total_attempts < 25:
        return K_FACTORS['intermediate']
    elif total_attempts < 50:
        return K_FACTORS['advanced']
    else:
        return K_FACTORS['expert']

def calculate_success_rate(
    successful_attempts: int,
    overall_attempts: int,
    correct_tests: int,
    performed_tests: int,
    time_used_ms: int,
    time_limit_ms: int = DEFAULT_TIME_LIMIT,
    discrimination: float = TIME_DISCRIMINATION
) -> float:
    if overall_attempts == 0 or performed_tests == 0:
        return 0.0
    attempt_ratio = successful_attempts / overall_attempts
    test_ratio = correct_tests / performed_tests
    time_component = max(0.0, discrimination * (time_limit_ms - time_used_ms))
    W = attempt_ratio * test_ratio * (1.0 + time_component)
    return max(0.0, min(2.0, W))

def update_elo_ratings(
    student_rating: float,
    question_difficulty: float,
    success_rate: float,
    k_factor: int
) -> tuple[float, float]:
    expected_score = calculate_expected_score(student_rating, question_difficulty)
    new_student_rating = student_rating + k_factor * (success_rate - expected_score)
    new_question_difficulty = question_difficulty + k_factor * (expected_score - success_rate)
    new_student_rating = max(RATING_MIN, min(RATING_MAX, new_student_rating))
    new_question_difficulty = max(RATING_MIN, min(RATING_MAX, new_question_difficulty))
    return new_student_rating, new_question_difficulty

async def detect_stagnation(user_id: str, current_module_id: str, db: AsyncSession) -> bool:
    if current_module_id == "CH03": return False
    result = await db.execute(
        select(AssessmentLog)
        .where(AssessmentLog.user_id == user_id)
        .where(AssessmentLog.is_final_attempt == True)
        .order_by(AssessmentLog.timestamp.desc()).limit(5)
    )
    last_5_logs = result.scalars().all()
    if len(last_5_logs) < 5: return False
    deltas = [log.theta_after - log.theta_before for log in last_5_logs 
              if log.theta_before is not None and log.theta_after is not None]
    if len(deltas) < 5: return False
    return numpy.var(deltas) < 165

def check_fallback_trigger(group_assignment, current_module_id, is_next_module_unlocked, recent_logs) -> bool:
    if group_assignment != 'A' or current_module_id == "CH03" or is_next_module_unlocked:
        return False
    if len(recent_logs) < 9: return False
    wrong_count = sum(1 for log in recent_logs if not log.is_correct)
    return wrong_count > 4.5
\end{lstlisting}


\subsection{Algoritma Re-ranking (Peer Matching)}
Kode implementasi untuk \textit{re-ranking} dapat ditinjau pada repositori dalam berkas bernama \texttt{peer\_matching.py} di direktori \texttt{app/core/}. Berikut adalah kode lengkapnya:

\begin{lstlisting}[language=Python, caption={Implementasi Peer Matching (Re-ranking)}, label={lst:peer_matching}]
import numpy
import random
from typing import Optional
from sqlalchemy.ext.asyncio import AsyncSession
from sqlalchemy.future import select
from sqlalchemy import func as sql_func
from app.db.models import User, PeerSession

COHEN_D_THRESHOLD = 0.5
FALLBACK_POPULATION_STD = 100.0

async def calculate_population_std(db: AsyncSession) -> float:
    result = await db.execute(select(User.theta_individu).where(User.status == "ACTIVE"))
    thetas = [row[0] for row in result.all()]
    if len(thetas) < 2: return FALLBACK_POPULATION_STD
    population_std = numpy.std(thetas, ddof=0)
    return float(population_std) if population_std != 0 else FALLBACK_POPULATION_STD

async def find_heterogeneous_peer(requester: User, db: AsyncSession) -> Optional[User]:
    population_std = await calculate_population_std(db)
    min_difference = COHEN_D_THRESHOLD * population_std
    result = await db.execute(
        select(User).where(
            User.user_id != requester.user_id,
            User.group_assignment != 'B',
            User.status != "NEEDS_PEER_REVIEW",
            sql_func.abs(User.theta_individu - requester.theta_individu) >= min_difference
        ).order_by(User.theta_individu.desc()).limit(5)
    )
    candidates = result.scalars().all()
    if not candidates: return None
    return random.choice(list(candidates))

async def create_peer_session(requester, reviewer, question_id, requester_query, db) -> PeerSession:
    peer_session = PeerSession(
        requester_id=requester.user_id,
        reviewer_id=reviewer.user_id,
        question_id=question_id,
        requester_query=requester_query,
        status="PENDING_REVIEW"
    )
    db.add(peer_session)
    requester.status = "NEEDS_PEER_REVIEW"
    await db.flush()
    return peer_session
\end{lstlisting}

